    \draw[flowarrow_red] (L3) -- (LTotal.south);
    \draw[flowarrow_red] (LT) to[bend right=15] (LTotal.south east);
    
\end{tikzpicture}
\end{center}
\end{frame}

% ---------------- SLIDE 19 ----------------
\begin{frame}{Backpropagation Through Time (BPTT)}
\small
Có 3 tham số ta cần phải tìm là $W, U, V$. Để thực hiện gradient descent, ta cần tính: $\frac{\partial L}{\partial U}$, $\frac{\partial L}{\partial V}$, $\frac{\partial L}{\partial W}$.

Tính đạo hàm với $V$ thì khá đơn giản vì nó chỉ phụ thuộc trực tiếp vào bước hiện tại:
\[ \frac{\partial L}{\partial V} = \sum_{t=1}^{T} \frac{\partial L_t}{\partial \hat{y}_t} \cdot \frac{\partial \hat{y}_t}{\partial V} \]

Tuy nhiên với $U, W$ thì lại khác. Xét tổng quát tại thời điểm $T$:
\[ \frac{\partial L}{\partial W} = \sum_{t=1}^{T} \frac{\partial L_t}{\partial \hat{y}_t} \cdot \frac{\partial \hat{y}_t}{\partial s_t} \cdot \frac{\partial s_t}{\partial W} \]

Do trạng thái $s_t = \tanh(W \cdot s_{t-1} + U \cdot x_t)$ có $s_{t-1}$ phụ thuộc vào $W$, ta phải áp dụng Chain Rule truy hồi liên tục về tận thời điểm đầu tiên.
\end{frame}

% ---------------- SLIDE 20 ----------------
\begin{frame}{BPTT: Đạo hàm cho W}
\small
Sử dụng công thức đạo hàm hàm hợp $(f(x) \cdot g(x))' = f'(x)g(x) + f(x)g'(x)$, ta có:

\[ \frac{\partial s_T}{\partial W} = \frac{\partial s'_T}{\partial W} + \frac{\partial s_T}{\partial s_{T-1}} \cdot \frac{\partial s_{T-1}}{\partial W} \]
trong đó $\frac{\partial s'_T}{\partial W}$ là đạo hàm của $s_T$ với $W$ khi coi $s_{T-1}$ là hằng số đối với $W$.

Tương tự, trong biểu thức $s_{T-1}$ có $s_{T-2}$ phụ thuộc vào $W$ \dots nên áp dụng công thức trên và chuỗi Chain Rule, ta được:
\[ \frac{\partial L}{\partial W} = \sum_{t=1}^{T} \sum_{i=0}^{t} \frac{\partial L_t}{\partial \hat{y}_t} \cdot \frac{\partial \hat{y}_t}{\partial s_t} \cdot \frac{\partial s_t}{\partial s_i} \cdot \frac{\partial s'_i}{\partial W} \]

trong đó $\frac{\partial s_t}{\partial s_i} = \prod_{j=i}^{t-1} \frac{\partial s_{j+1}}{\partial s_j}$ và $\frac{\partial s'_i}{\partial W}$ là đạo hàm của $s_i$ với $W$ khi coi $s_{i-1}$ là hằng số.
\end{frame}

% ---------------- SLIDE 21 ----------------
\begin{frame}{BPTT: Đạo hàm cho U}
\small
Hoàn toàn tương tự với ma trận trọng số $U$ (từ input vào hidden layer):
\[ \frac{\partial L}{\partial U} = \sum_{t=1}^{T} \sum_{i=0}^{t} \frac{\partial L_t}{\partial \hat{y}_t} \cdot \frac{\partial \hat{y}_t}{\partial s_t} \cdot \frac{\partial s_t}{\partial s_i} \cdot \frac{\partial s'_i}{\partial U} \]

\vspace{0.4cm}
\begin{alertblock}{\textbf{Cốt lõi vấn đề của BPTT}}
Nhìn vào công thức tính đạo hàm của $L$ với $W$ và $U$ ở trên, ta thấy xuất hiện nhân tử:
\[ \frac{\partial s_t}{\partial s_i} = \prod_{j=i}^{t-1} \frac{\partial s_{j+1}}{\partial s_j} \]
Đây chính là nguồn cơn của hiện tượng \textbf{Vanishing Gradient} (Triệt tiêu đạo hàm) hoặc \textbf{Exploding Gradient} (Nổ đạo hàm) khi chuỗi $T$ lớn.
\end{alertblock}
\end{frame}

% ---------------- SLIDE 22 ----------------
\begin{frame}{Chứng minh Toán học: Phân rã Gradient trong RNN}
\footnotesize
Xét thành phần cốt lõi trong BPTT là đạo hàm trạng thái qua nhiều bước thời gian: $\frac{\partial s_T}{\partial s_k} = \prod_{j=k}^{T-1} \frac{\partial s_{j+1}}{\partial s_j}$.\\[0.1cm]
Do $s_{j+1} = \tanh(W \cdot s_j + U \cdot x_{j+1})$, ma trận Jacobian tại mỗi bước là: $\frac{\partial s_{j+1}}{\partial s_j} = \text{diag}(1 - s_{j+1}^2) \cdot W$.

\vspace{0.05cm}
\begin{block}{\textbf{Đánh giá Chuẩn Ma trận (Matrix Norm)}}
Áp dụng tính chất $\Vert A \cdot B \Vert \le \Vert A \Vert \cdot \Vert B \Vert$, ta lấy chuẩn 2 vế:
\[ \left\Vert \frac{\partial s_{j+1}}{\partial s_j} \right\Vert \le \Vert \text{diag}(1 - s_{j+1}^2) \Vert \cdot \Vert W \Vert \le 1 \cdot \Vert W \Vert = \Vert W \Vert \]
(Vì đạo hàm của $\tanh$ lớn nhất bằng 1 nên $\Vert \text{diag}(1 - s_{j+1}^2) \Vert \le 1$).
\end{block}

\vspace{0.05cm}
\begin{alertblock}{\textbf{Cơ chế Vanishing Gradient}}
Nếu trị riêng lớn nhất của $W$ nhỏ hơn 1 (tức $\Vert W \Vert = \gamma < 1$): $\left\Vert \frac{\partial s_T}{\partial s_k} \right\Vert \le \gamma^{T-k} \xrightarrow{T-k \to \infty} 0$.\\[0.1cm]
$\Rightarrow$ Tín hiệu gradient suy giảm theo hàm mũ về 0. Trọng số mất hoàn toàn khả năng cập nhật theo các lỗi từ quá khứ xa.
\end{alertblock}
\end{frame}

% ---------------- SLIDE 23 ----------------
\begin{frame}{Hiện tượng Nổ Gradient (Exploding Gradient)}
\small
\begin{alertblock}{\textbf{Cơ chế Toán học (Exploding Gradient)}}
Ngược lại với triệt tiêu, nếu trị riêng lớn nhất của $W$ lớn hơn 1 ($\Vert W \Vert = \gamma > 1$) và phần lớn nơ-ron hoạt động ở vùng tuyến tính của $\tanh$ (nơi đạo hàm xấp xỉ 1):
\[ \left\Vert \frac{\partial s_T}{\partial s_k} \right\Vert \approx \gamma^{T-k} \xrightarrow{T-k \to \infty} \infty \]
\end{alertblock}

\vspace{0.1cm}
\begin{block}{\textbf{Hậu quả Hình học ("Vách đá" Hàm Loss)}}
Sự nhân dồn liên tục của $W$ qua $T$ bước thời gian tạo ra những "vách đá" (cliffs) cực dốc trong không gian tham số. 
\begin{itemize}
    \item Khi mô hình tiến gần vách đá, đạo hàm bùng nổ theo hàm mũ.
    \item Gradient khổng lồ gây ra bước cập nhật (Overshooting) đẩy văng mô hình khỏi vùng hội tụ, sinh ra lỗi \texttt{NaN} hoặc \texttt{Inf}.
\end{itemize}
\end{block}
\end{frame}

% ---------------- SLIDE 24 ----------------
\begin{frame}{Giải pháp khắc phục: Gradient Norm Clipping}
\footnotesize
Để ngăn overshooting, ta hãm phanh bước nhảy. Đặt vector $g = \nabla L = \left[ \frac{\partial L}{\partial W}, \frac{\partial L}{\partial U}, \frac{\partial L}{\partial V} \right]$ là vector gradient chứa đạo hàm của tất cả các trọng số trong mô hình.

\vspace{0.05cm}
\begin{exampleblock}{\textbf{Cơ chế Cắt tỉa (Clipping)}}
Nếu độ lớn (norm) của vector $g$ vượt ngưỡng $\theta$, ta co giãn toàn bộ vector:
\[ g_{\text{clipped}} = \begin{cases} 
g & \text{nếu } \Vert g \Vert \le \theta \\
\theta \cdot \frac{g}{\Vert g \Vert} & \text{nếu } \Vert g \Vert > \theta 
\end{cases} \]
\end{exampleblock}

\vspace{0.05cm}
\begin{alertblock}{\textbf{Tại sao dùng Norm Clipping thay vì Value Clipping?}}
\begin{itemize}
    \item \textbf{Value Clipping (Cắt từng phần tử):} Ép lẻ từng $g_i \in [-\theta, \theta]$. Việc này \textbf{làm bóp méo hướng} của gradient, khiến mô hình cập nhật sai hướng.
    \item \textbf{Norm Clipping (Co giãn vector):} Phép nhân vô hướng với tỷ lệ $\frac{\theta}{\Vert g \Vert}$ giúp vector mới tuy ngắn lại nhưng \textbf{bảo toàn tuyệt đối hướng đi nguyên thủy} trong không gian đa chiều.
\end{itemize}
\end{alertblock}

\vspace{0.05cm}
\textit{\textcolor{crimsonred}{*Lưu ý:}} Clipping chỉ chống Nổ đạo hàm, hoàn toàn vô dụng với Triệt tiêu đạo hàm.
\end{frame}



% ---------------- SLIDE 25 ----------------
